Financial analysts reviewing bank results need to understand not only the reported figures but also what those figures represent. Similar measures can differ in accounting basis, reporting period or operational scope, making it important to check the original evidence before using an answer in a financial briefing.

FinTrace is a retrieval augmented generation (RAG) prototype that helps analysts explore six official Commonwealth Bank of Australia and National Australia Bank reports covering FY2024 and FY2025. Users can find financial figures, compare results across years within a bank, understand financial terms and explore management's explanations. Answers include relevant financial context and references to the original report pages. Python performs supported calculations using checked source values, while a local language model helps select relevant passages and explain financial terms.

Development followed an iterative testing process covering document extraction, chunking, retrieval and answer production. Testing revealed problems including figures separated from their labels, incorrect associations between financial measures and incomplete answers. These findings guided improvements to evidence handling and response coverage. Investigating one apparent system error also uncovered a mistake in the expected answer, demonstrating the importance of checking evaluation references against the original reports.

The resulting prototype brings financial answers, supporting evidence and calculation workings into one interface. It supports an analyst's review while making evidence limitations visible. Some explanations remain incomplete and response times vary, providing clear priorities for further development. FinTrace demonstrates how testing throughout the RAG pipeline can produce a more useful and accountable tool for exploring financial reports.
